% ============================================================
%  H. Høffding, "Forholdet mellem Tro og Viden i dets
%  historiske Udvikling. Et indledende Foredrag."
%  Først trykt 1885; optrykt i: Mindre Arbejder,
%  Kjøbenhavn: Gyldendal, 1899.
%  TRANSLATION (English) --- Hans Halvorson
%  COMPLETE: Sections I--VI translated from transcription.tex.
%
%  Editorial notes on the translation:
%   - The Dante quotation (Inferno I.124--126) is given in the Danish verse
%     rendering the printed text uses; it is Englished as verse, not replaced
%     by an existing English Dante translation.
%   - The Goethe couplet ("Wer immer strebend sich bemüht") is left in German,
%     as printed.
%   - Danish work-titles in the author's references are left in Danish
%     (Martensen's "Kristelige Etik", "Dogmatik"); German titles likewise
%     (Lange's "Geschichte des Materialismus").
%   - Schleiermacher's title is given in English ("Speeches on Religion to its
%     Cultured Despisers"), matching the Danish text's own translated form.
%   - Danish »...« → English curly doubles throughout.
%   - The piece ends, as the print does, on a dash after section VI.3.
% ============================================================
\documentclass[12pt,a4paper]{article}
\usepackage[utf8]{inputenc}
\usepackage[T1]{fontenc}
\usepackage{libertinus}
\usepackage{libertinust1math}
\usepackage[english]{babel}
\usepackage[protrusion=true,expansion=true]{microtype}
\usepackage{setspace}
\usepackage[margin=1.2in]{geometry}
\usepackage{fancyhdr}
\usepackage{hyperref}

\hypersetup{
  pdftitle={The Relation between Faith and Knowledge in its Historical Development},
  pdfauthor={H. Høffding},
  hidelinks
}

\pagestyle{fancy}
\fancyhf{}
\rhead{Høffding, \emph{Faith and Knowledge in its Historical Development} (1885)}
\cfoot{\thepage}

\setstretch{1.3}

\title{\textbf{The Relation between Faith and Knowledge\\[4pt]
       in its Historical Development}\\[6pt]
  \large An Introductory Lecture}
\author{Harald Høffding\\[4pt]
  \small Translated from the Danish by Hans Halvorson}
\date{Printed in: \emph{Mindre Arbejder}, Copenhagen: Gyldendal, 1899\\
  (Lecture delivered 1885)}

\begin{document}
\maketitle
\thispagestyle{fancy}

\bigskip

\noindent The problem of faith and knowledge arises of necessity at a
certain stage of cultural development. We encounter it in the Indian,
the Greek, and the Muhammadan worlds as soon as intellectual and
spiritual development has reached a certain point. Here, however, we
shall confine ourselves to tracing the main forms in which it has
appeared within the European spiritual development shaped by
Christianity.

\section*{I.}

No instruction on this question can be drawn from the source writings
of Christianity. On the standpoint from which they were composed, no
knowledge distinct from faith exists. The gaze is fixed on one single
object, one single goal, which lays claim to all striving and all
spiritual force. The endeavor to know the truth is here immediately one
with the urge to win the salvation of the soul, and does not go its own
way by its own method. There is no distinction between theory and
practice, since theory and theoretical interest simply do not exist.
When the Deity itself has revealed itself visibly to all, --- when faith
makes it possible to attain understanding of the Mystery that contains
all the treasures of wisdom and knowledge, --- every motive for seeking
knowledge by other paths is lacking.

It was natural that faith should from the outset have to be all in all,
leaving no room for other spiritual needs. The first generation still
stood so close to the mighty spiritual movement that led to the founding
of the Christian community that it lived and breathed entirely in the
great memories and visions. Added to this was the fact that this first
generation lived in the belief in the imminent return of Christ. Then
everything was to be concluded, the last judgment rendered, and all human
and earthly arrangements give way to a higher order of things. Who, with
such a catastrophe in view, had the time and composure of mind for
philosophical thinking or scientific inquiry? Science could no more than
any other human endeavor acquire any significance. There was to be no
historical development at all; the goal had been reached, and the end
stood at the door. One does not understand the New Testament and its
relation to cultural endeavor in its various forms unless one holds fast
to this. What happened here with scientific striving was the same as with
the striving for personal freedom and with the entering into marriage
(see 1~Corinthians, ch.~7). How could all such things acquire
significance --- or even become the object of a drive --- when all
thoughts and energies had to be strained to be prepared for the great,
the absolute decision!

\section*{II.}

Time passed; the conclusion did not come. The Church consoled itself with
the thought that with God a thousand years are as one day, and one day as a
thousand years (2~Peter, ch.~3), and soon began to settle down in the world,
to adapt itself to circumstances, and to take up a position toward the
various sides of the endeavors of cultural history. Toward science its
position became this: that it no longer found in faith all requisite
knowledge, but sought to show that the teachings of faith set the crown upon
all human knowledge. The Church appropriated the science the Greeks had
developed --- Greek astronomy, medicine, and philosophy --- and placed its
theology as the spire upon the edifice of science, much as it converted the
Greek temples into churches and consecrated, for example, the Parthenon,
Athena's temple, to ``holy Wisdom.''

Already in the so-called Gospel of John we find, in the opinion of some,
Platonic philosophy and Christian faith joined together. This endeavor is
carried on by the Church Fathers and reaches its completion in the classical
period of theology among the great teachers of the medieval Church. In
Thomas Aquinas (born 1227, died 1274), the teacher still honored and
acknowledged in our own day by the Catholic Church, the relation between
faith and science stands as follows. Science investigates all subordinate and
particular causes, while the thought of the believer concerns the first, the
absolute cause of all things. Only theology, therefore, teaches us rightly to
understand things; science really teaches us only what things are, not whence
they derive. But there is no conflict between faith, which is the principle of
theology, and reason, which is the principle of science. Reason stems from God
just as much as faith does, and faith does not conflict with reason, although
it goes beyond reason.

In poetic form we find this characteristic medieval standpoint in Dante (born
1265, died 1321), who was a disciple of Thomas. In his \emph{Commedia} the
poet relates how, in order that his soul might be purified, it was granted him
to pass through the world beyond: Hell, Purgatory, and Paradise. It is his
transfigured beloved, Beatrice --- the symbol of theology grounded in heavenly
revelation --- who bestows this favor upon him and sends Virgil, the symbol of
natural science, to lead him through the underworld. But Heaven Virgil dares
not tread; he himself states the reason, with great meekness:

\begin{quote}
for He who rules the city of the sky\\
wills not that I should cross its threshold, seeing\\
that I against His laws a rebel stood.
\end{quote}

Therefore, when the wanderers draw near to Heaven, Beatrice reveals herself,
and Virgil disappears.

So pious and so humble had natural reason become in the Middle Ages. But its
powers grew, and the close of the Middle Ages is marked by violent discord
between Beatrice and Virgil, faith and knowledge. Here we meet the problem
posed for the first time.

\section*{III.}

The proud edifice that Thomas Aquinas had completed, and that Dante had sung,
soon began to totter on its foundation. Continued reflection discovered that
it had been built of elements that could not be brought into harmony. Already
William of Occam (\dag\ 1347) finds a contradiction between natural reason and
the divine omnipotence that is exalted above every law. He teaches that
something may be true in theology which is not true in philosophy: if the
Bible and the Church teach something that is self-contradictory, it must
nevertheless be taken for truth! Theology rests on supernatural authority and
is therefore no science.

Have we then two kinds of knowledge, two kinds of logic? This question, which
lies so near at hand, is not brought forward. Occam himself was a believing
Catholic; he pointed out the discord, but did not pursue its consequences. The
Catholic Church held fast to its holy Thomas and has permitted no other view
of faith and knowledge to assert itself within its precincts. In our own
century no fewer than three Catholic philosophers have been condemned, because
they were either more or less rationalistic than St.~Thomas. For the Catholic
Church there no longer exists any problem here. --- Let us see, then, how the
matter stands within Protestantism.

\section*{IV.}

In modern times life spreads out in a great many different directions and has
a far more variegated and many-sided character than in antiquity and the
Middle Ages. And the Church stands as only one among the many domains of the
movement of life; it is no longer all in all, as in the first community, nor
the crown of the whole, as in the Middle Ages. History is not first and
foremost a history of the Church, but, among other things, also a history of
the Church. It is predominantly passively that the Church receives the impact
of the new directions of life and seeks to come to terms with them; but they
no longer proceed from it.

The relation of ecclesiastical Protestantism to science is much the same as
Catholicism's. Luther had indeed appealed --- not only to Holy Scripture, but
also to ``clear reasons''; but otherwise he applied to reason Paul's words,
that the woman shall keep silence in the congregation. Reason, then, was to
have Virgil's role, and might not enter the Holy of Holies. The Protestant
Church would gladly have gone on performing ``the divine comedy.'' But Virgil
has proved no longer tractable.

Protestant orthodoxy became a new scholasticism, but more anxious and less
grand than the medieval. --- It has found no Dante. The great Protestant poets
pass it by or go against it. Poetry's relation to the Church forms a parallel
to that of science: the ancient Church has no poetry, Catholicism an
ecclesiastical poetry, Protestantism a poetry standing outside the Church, if
not an anti-ecclesiastical one.

For Protestantism the eighteenth century is what the fourteenth was for
Catholicism. But the crisis it brought was far more thoroughgoing, and reached
far deeper, than the Renaissance movement. It is the birth-time of the
independent humane view of life. Forerunners there were, of course, many ---
but only now were they understood and acknowledged. Only now did Spinoza, for
example, who for a century and a half had lain ``like a dead dog,'' find a
circle able to think his thoughts and make them their own. Only now did the
theological view of life acquire an opponent of equal standing. I am thinking
here not of any definite, finished system, but of a general direction of
thought and life, which in its broad features was common to such men as
Lessing and Kant, Schiller and Goethe. --- As one sees, this new view of life
found at once more than one Dante.

What is distinctive about this standpoint --- whose leading thoughts were only
transiently cast into shadow by the Romantic reaction --- is, seen from the
theoretical side, that what it sets over against theological faith is not a
pretended insight into the impossibility of the theological doctrines, but an
assertion of their unprovability. Kant's great achievement was, seen from this
side, to give a scientific doctrine of the limits of science. Not humbled and
cowed, like a rebel who has submitted, not checked from without, but with
independent and clear self-knowledge, human reason here says: so far can I
reach, and no farther! --- The limits of knowledge follow from the nature of
knowledge. By an investigation of the ultimate presuppositions and laws of our
knowledge, the domain within which it can move is staked out. This domain is
beyond surveying, yet it is definitely bounded over against that which by its
nature cannot become an object of our knowledge, constituted as this knowledge
once and for all is.

For Kant and his age the significance of knowledge lay not in its giving us
``the absolute truth,'' but in this, that by means of it we can advance
continually in truth. Eternal seeking, eternal striving became from now on the
watchword:

\begin{quote}
\itshape
Wer immer strebend sich bemüht,\\
den können wir erlösen.
\end{quote}

Not as though nothing could be found or attained: but every find, every goal
attained, stands as the starting point for new striving. It was a great law of
life and of the world that was here given expression --- the thought of
development in one of its most important forms, a thought to which our own age
has given a realistic underpinning and a closer critical determination, but
which it is that generation's honor to have uttered for the first time with
full and clear consciousness.

But still more important was another principle, which was laid down as the
main point in the humane view of life: the principle of the free conscience.
Kant's second great achievement was that he maintained the possibility and the
necessity of a purely humane ethics, and established that conscience can as
little as theoretical inquiry have its limits staked out for it from without.
The measure of all that we are to acknowledge as good, noble, and exalted must
be found within ourselves and be determined by human nature.

When those who stand at this standpoint cannot accept ecclesiastical faith, it
is because it conflicts not only with their understanding, but with their
conscience. One is not merely to assent to a formulation of the truth given
once and for all, but also to bow before doctrines that cannot be brought into
agreement with conscience's measure of good and evil.

Only here does the problem become sharpened, and its real sting come to light.
What from the one side is made the highest duty, in a way the only duty, a
duty whose non-fulfillment brings damnation --- that is rejected from the
other side precisely as conflicting with the clear and definite demand of
conscience! ---

Bishop Martensen saw clearly, from his own standpoint, that the most important
turning point in the controversy over faith and knowledge came with the
founding of ethical humanism at the close of the last century. In his
\emph{Kristelige Etik} he deals at length with the representatives of this
tendency and their relation to the faith of the Church. The final explanation
of their refusal to accept it he finds, rightly, not in their understanding,
but in their will --- only that this will is, according to his assertion,
defiance against the Holy: ``Against God's holiness they have a secret
antipathy!''

How strange, or rather how appalling, human life must look from the orthodox
standpoint! The most honest seeking and striving, the strictest
conscientiousness, joined with the most unreserved acknowledgment of the
imperfection of what has been attained, must from this standpoint be construed
as antipathy toward holiness. --- If one has not understood why the problem of
faith and knowledge so easily sets the passions in motion, one understands it
here.

\section*{V.}

The new humane standpoint seemed to have a fairly good prospect of prevailing.
The most eminent minds of the age paid homage to it, the leading statesmen
worked in the direction it marked out, within the Church itself orthodoxy was
decidedly weakened, and confessional differences receded before an enlightened
and liberal-minded way of thinking that was common to all. Such were
conditions at the beginning of the century, not only in the Protestant, but
also in the Catholic Church.

But how differently were conditions to develop in the course of the century,
and how changed were they already after the first two decades had run out!

A single trait shows the change plainly. In the year 1799 Schleiermacher
published his \emph{Speeches on Religion to its Cultured Despisers}. He was by
no means contending there for any dogmatic faith, but asserted a pantheistic
religiosity and glorified Spinoza. For such a view he wished to win
recognition among the cultured. But when in the year 1821 he saw the third
edition of his book through the press, he declared in the preface that if one
now looked about among the cultured, one would rather find it needful to write
speeches to bigots and slaves of the letter, to ignorant and uncharitable
condemners, to the superstitious and the hyper-dogmatic!

Four years earlier (1817) the Holstein pastor Claus Harms had, in commemoration
of Luther's action, published 95 new theses, in which he protests, among other
things, against ``the popes of our time'': reason in the theoretical domain,
and conscience in the practical domain. --- How was it possible that such a
performance could win adherence and mark the beginning of a new period?

What, on the whole, was the reason that the breakthrough of humanism at the
close of the eighteenth century did not come to inaugurate a development
carried on in a straight line down to the present, but was succeeded by its
sharp and conscious opposite?

Since a closer psychological-historical explanation of the mighty reversal
that came to fill so large a part of our century, and made it ``the age of
oppositions,'' may serve to throw further light on the nature of our problem,
we shall enter into it somewhat more closely, but in such a way as to keep
chiefly to the causes that can be found in the very character and
distinctiveness of the humanism of the eighteenth century. Many different
causes were of course at work together. Thus the general political reaction
had its influence on religious development as well; the resurrection of the
principle of authority in the state had also to lead to its being favored in
the religious domain, and the political reactionaries often made use of
religion in the most repellent way. But had there not been deep-lying inner
causes, this outward influence could not have accomplished so much. ---

\medskip
\noindent 1. Every reaction is due in large part to an effect of contrast.
Just as the eye wearied by light finds rest in darkness, or at least in a
fainter gleam, so on the whole the human spirit cannot endure a strong and
protracted straining of a single side of its nature, but throws itself
involuntarily over to the opposite side in order to find a new position of
rest. With respect to the problem we are speaking of here, it is above all the
opposition between the active and the passive, the effecting and the receptive
side of our nature that becomes significant. Not as if in our spiritual life
there could be pointed out states or phenomena in which we are either
absolutely effecting or absolutely receptive: in every form of spiritual life
and in every spiritual state there can, on the contrary, be pointed out both an
activity and a passivity. But these two moments can stand in very different
relations to one another; now the one, now the other has the preponderance.
Thus inquiring and brooding thought, the testing and forward-striving
conscience, and the forming and creating imagination have a preponderantly
active character, and over against them in this respect stand the need for
consolation and self-surrender, and feelings such as piety and resignation.
Where it is a matter of bringing forth new forms of life, the active side of
our nature acquires the preponderance and lays claim to the whole of our
energy. Just such a task presented itself to the generation of the latter part
of the eighteenth century. It was the age in which unlimited advance was the
watchword, and inertia was regarded as original sin. But there had to come a
time when the other sides of human nature demanded their satisfaction, and
when men sought rest after the restless labor. And rest was not to be found in
the new forms. They still required all too much active cooperation, made a
constant demand for self-activity. Rest, on the other hand, might perhaps be
found in the old forms, on which in one's onward rush one had turned one's
back, but of which one now discovered how wonderfully they were adapted to the
needs of the life of feeling: for centuries or millennia generation after
generation had deposited its experiences in them, had arranged them so that
uninterrupted labor was not needed to use them. Our activity seeks the new;
with our effecting nature we live in the present; but the passive and
receptive side in us is determined and nourished preponderantly by the
after-effects of the preceding age; here spiritual life draws near to the
unconscious and takes nourishment from it. An example of this from the life of
the individual we have in the power of childhood memories: they work without
our conscious cooperation, and they rise up with overwhelming force
particularly when the mind is weary of active intellectual labor.

It is characteristic that Schleiermacher, precisely at the transition to the
new century, set up his famous definition of religion as the feeling of
unconditional dependence. It was like a program for the direction spiritual
life was now to take. Humanism was succeeded by confessionalism, free inquiry
and the free conscience by belief in authority.

Are we then consigned to a labor of Sisyphus? Does the stone always roll back
once we have got it up the mountain? --- No; but what lies in the foregoing is
that only that truth can be victorious which is able to penetrate and fill the
whole of life. And that does not happen at a single stroke; it takes time, and
only through many oscillations does it advance. It seems to be a psychological
law that knowledge develops more rapidly than the life of feeling. Only slowly
are the results of active and conscious labor deposited in the involuntary life
of feeling. Before that can happen there naturally comes a period in which the
new thoughts grow pale and prove unsatisfying; doubt and unrest take hold of
many, and there seems to be no other way out than to throw oneself back into
the old forms of the life of feeling, if one is not to perish. But this
disharmony will be only transitory: if the new thoughts are valid, if they
really express some of the laws to which our existence is subject, then they
will always work their way forward anew, and the time of unstable equilibrium
will be past. In the long run the life of feeling will develop in greater or
lesser harmony with the life of thought: it was in that way, after all, that
those very old forms and dogmas came into being. Even the periods of reaction
themselves prove to bear the stamp of what they combat; into the old
wineskins new wine will to no small extent be poured, and only one who keeps
to externals will suppose that everything has become again as it was before.

We stand here before a spiritual connection of nature to which we must bow,
and which we cannot defy. Whoever has come to see it rightly will neither
share the cheaply bought confidence of many in the nearness of victory, nor
lose heart at an apparent setback. He knows that just as the apparently sudden
revolutions signify only that what has long been prepared in secret now steps
forth into the light of day, so the sudden reactions indicate that there were
drives and needs which were only imperfectly satisfied by the seemingly so
brilliant advance, and which therefore must seek satisfaction in the old,
until they have accommodated themselves to the new, or the new to them.

\medskip
\noindent 2. The humanism of the eighteenth century not only gave the knowing
and effecting part of our nature a decided preponderance over the feeling and
receptive part, but, in close connection with this, it had a decidedly
intellectually aristocratic character. The whole movement was, and by its
nature had to be, confined to a small circle; the people proper were very
little affected by it. At no time, perhaps, have the well-placed classes in
society shown so great a zeal for taking the lead in the work of enlightenment
and progress. But the distance between those who led and those who needed
leading was extraordinarily great. And the development had led to a break with
the common faith, the common forms for the highest thoughts. In the Middle
Ages the most deeply brooding schoolman and the simplest serf could meet in
the same confession. Now this unity was shattered, and there had to arise ---
especially since at the same time feeling for the national and the popular took
a strong upswing --- a need to be one with one's people in faith as well. The
reaction against the eighteenth century acquired, seen from this side, an
intellectually democratic character. It was, after all --- to take a single
example --- the conviction that ``the faith of the unlearned cannot possibly be
bound to the testimony of the learned'' that led Grundtvig to his
``discovery.''

Like the discord between the different sides or forces in the spiritual life,
the discord within the life of the people will also be a constantly operative
motive for returning to the old faith. Common forms and expressions for the
highest thoughts are created slowly and presuppose that the thoughts have
passed into the flesh and blood of all. Since now the individual differences
in religious need and religious sense must be assumed to rise and increase the
more religious freedom becomes a reality, such a community will become ever
harder to bring about.

A testimony to how great the need for community can be is offered by the
philosophy of religion of Schelling and Hegel. These men and their followers
were convinced that in their philosophical ideas they had developed the same
content that the dogmas of the Church expressed. The difference was to concern
only the form: what theology has as representation, in the form of imagination
and feeling, philosophy was to shape into abstract concept, without the
essential being thereby lost. In the dogma of God's becoming man, for example,
there was supposed to be contained the philosophical thought that even the
highest must develop through struggle and oppositions, and that a principle
shows its validity only when it is able to penetrate actual life with all its
discord and all its conflicting relations.

In and for itself it is entirely legitimate to regard the dogmas as symbols.
For the non-believer they cannot be anything else, if they are to be anything
at all. And many who consider themselves very believing, and are so considered
by others, stand in reality at this standpoint. But it will always lead to
dubious consequences when one uses the symbolic interpretation to bring about
a false semblance of community. A real unity cannot be produced by that road.
On such a false semblance rested the last attempt at a reconciliation of faith
and knowledge that has been made, Hegel's philosophy of religion, from which
many expected a ``new era.'' Here Virgil was not to subordinate himself humbly
to Beatrice; they were to walk arm in arm. But this apparent agreement gave
occasion only for the opposition between faith and knowledge to be expressed
in a still far sharper and more pointed form than hitherto, as was done by
Strauss and Ludwig Feuerbach in Germany and by Søren Kierkegaard among us.
These men were agreed that the content of theological faith was the
incomprehensible; the difference between them consisted in this, that Strauss
and Feuerbach said: It is incomprehensible, and therefore it cannot be
believed! while Kierkegaard said: Precisely because it is incomprehensible, it
can and shall be believed! --- Strauss's words: ``Of false attempts at union
enough has been done; only a sharpening of the oppositions can lead further!''
are characteristic of the course of development in the nineteenth century. A
famous scholar of church history, F.~C.~Baur, has rightly found the
distinctive mark of our century in this, that all oppositions --- in science,
religion, and politics --- are worked out to the highest possible degree. That
hangs together naturally with the fact that reaction and revolution in the
various domains have so frequently stood over against one another; thereby
there has arisen a clear consciousness of the nature of the oppositions and an
endeavor to work each of them out separately in its full consequence. Both the
necessity of making a choice between irreconcilable oppositions, and the
character of the paths one has to choose between, dawn more and more upon the
general consciousness. ---

If the mere need for community has thus led to sympathy with the old faith ---
whether this sympathy led to personal adherence to its letter or merely to a
symbolic interpretation of its statements --- then this whole question of
faith and knowledge had to present itself in a particularly serious way to
those who saw its connection with the social question. Even granting that the
scientific criticism of the dogmas leads to their dissolution, the thought will
still arise that these dogmas, however contrary to reason they may be, are for
many their only spiritual support and consolation --- that which keeps up
their hope under life's heavy struggle --- the only thing at their disposal
that can lift their gaze beyond the narrow circle in which their life moves.
It is from this point of view that one of the noblest philosophical thinkers of
our time, Albert Lange, the author of \emph{Geschichte des Materialismus}, has
expressed himself against Strauss and Feuerbach, with whom he agrees as
regards the purely theoretical side of the matter. It is on the whole
characteristic how great a role the question of religion's usefulness and
indispensability plays in the religious debate. What is at issue in the first
line is not: truth or not truth? --- but: consolation or not consolation? ---
In the debates of earlier times this side of the matter did not come forward
as something on its own, still less as the main thing. Men did not search for
grounds and motives for holding to dogmatic faith in spite of the
difficulties; rather, it was those who gave up dogmatic faith upon whom it fell
to give reasons for their step.

It is surely of very great importance that spiritual community should not
disappear, and that our race should not lose any force that can hold it up in
the struggle for existence. It is with full right that these considerations
become decisive for the manner in which the religious struggle is conducted;
but the struggle itself they can damp only for a time, since they point out no
way to the solution of the problem.

\medskip
\noindent 3. Finally, the humane view of life, as the eighteenth century was
able to present it, had a one-sidedly abstract and idealistic character. It
was stated according to its general principle and with the most important
presuppositions and demands it carries with it, but it lacked sufficient
carrying-through and a real foundation. Practically every advance the
nineteenth century has made has contributed its share toward remedying these
defects to some degree. The magnificent upswing of the natural sciences has
given a firmer basis for the conviction of the law-ordered connection of
nature. In the concept of development a guiding idea has been indicated around
which all our knowledge in the domain of nature and of history can gather.
Historical criticism and comparative religious science have widened and
corrected the material necessary for the treatment of problems in the
philosophy of religion. A more deeply penetrating psychological understanding
of the various sides of human spiritual life in their relation to one another,
and a renewed discussion of epistemological questions, give the problem of
faith and knowledge a somewhat different character from what it could formerly
have. --- All this is of course only small progress in relation to the
magnitude of the problem. The problem will for long ages accompany the human
race on its way, and the different ways in which it is conceived and treated
will be able to furnish a measure of the degree and direction of spiritual
development at the various times. Here we have only to show how the problem
must present itself in our own day.

\section*{VI.}

In the discussion of faith and knowledge that took place in our literature in
the sixties, great weight was laid on the distinction between faith and
theology, and I myself, as a disciple of Rasmus Nielsen, belonged to the number
of those who attached no slight hopes to it. It plainly does not help us, and
in any case it is not tenable. As soon as faith has a definite content, a
definite object, there must also develop a doctrine about this content, which
seeks to exhibit its inner coherence and its significance. Otherwise neither
the believer himself nor anyone else knows what it really is that he believes
in.

Nowadays, moreover, theology has given up supporting its assertions with
scientific proofs. It has given up all proud dreams of getting Virgil to walk
in leading-strings, and is content when it is left unassailed from the side of
scientific criticism. It declares its dogmas to be unprovable, perhaps even
concedes that they present difficulties for the ``natural'' understanding, and
takes essentially the standpoint of maintaining dogmatic faith as a spiritual
necessity of life --- a necessity, be it noted, which it is a duty to have. It
is supposed to be a sign of hardening and of defiance against the Holy when one
does not feel that need which makes reason bow before the incomprehensible.
Therein lies, as I have already had occasion to touch upon, the real sting of
the problem.

On the other hand, no scientific thinker will assert that the absolute
impossibility of the content of the dogmas can be strictly proved. In
scientific inquiry we build upon the principle of natural causes. It has helped
us through many a time, and hardly any other scientific principle for the
understanding of existence will ever be able to be set up; --- but it cannot be
proved that everything in existence is subject to this principle.

Indeed, we must go a step further still. Simultaneously with the wonderful
advances that have been made in our knowledge of nature and of history, both as
regards particulars and as regards the great connection of things, we have
learned in a far deeper sense than before that the whole of existence is one
great riddle. For the philosopher the least phenomenon is in this respect just
as instructive as the greatest and most conspicuous. The simplest causal
relation gives us more than enough to think about --- really encloses the whole
riddle of the world within itself. Everything we understand and explain, we
understand and explain by exhibiting a causal connection: if we only knew what
causal connection really is --- what it is that holds everything in nature
together and makes it that wonderful whole which reveals itself to us more and
more --- then all riddles would be solved! But to such knowledge every prospect
is barred. Our sensations, representations, and concepts (and thus also the
concept of cause) are quite good aids for finding our way about in existence,
but we have no right to suppose that they can depict for us the innermost
essence of existence. We can at most form speculative hypotheses about how
existence might be constituted in itself; but such hypotheses cannot, in the
nature of the case, receive any confirmation from experience.

From this it by no means follows that one bows before every theological oracle.
It follows only that the discussion of faith and knowledge must be conducted in
a somewhat different way from that in which it has hitherto been conducted. Let
us briefly touch on the main questions that will here have to be discussed.

\medskip
\noindent 1. When theological faith acknowledges that its assertions are
unprovable, or perhaps even irreconcilable with some of the most important
presuppositions and results of science, it apparently places itself at an
inaccessible and unassailable standpoint. There seems no longer to be any
common ground on which the battle can be fought.

But when theological faith nevertheless --- in order to express its content ---
makes use of representations and concepts that natural reason and scientific
thinking also employ, then it must be asked whether these concepts are taken in
exactly the same sense. It has proved to be bound up with insuperable
difficulties, even for the boldest speculation, to carry through propositions
and judgments about what lies beyond every possible experience. What, now, can
it really be that enables faith to come off better here than scientific
thinking? Why do the great difficulties in the dogmas fall away when one passes
over from knowledge to faith? A creation out of nothing, for example --- a
first cause --- an absolutely good being who is nevertheless the cause of
everything, and therefore also of evil --- such concepts, against which human
thought has collided again and again, until at last it has given up the attempt
to reconcile them with its other concepts --- how can they appear in a wholly
different light because they are declared to be objects of faith and not of
knowledge?

If it is answered (as, for example, by Bishop Martensen in his
\emph{Dogmatik}) that the rebirth of human nature that takes place in the
believer embraces knowledge as well, so that what was formerly incomprehensible
or even self-contradictory becomes intelligible, then of course the unreborn
cannot have a say about how things look after so radical a change. But since
the believing and reborn do not, after all, entirely give up their connection
with the natural world, and, insofar as they continue to belong to it, must
also continue to use their natural reason, they ought nevertheless to be able to
compare the two kinds of knowledge, the natural and the reborn, and tell us
wherein the difference really consists. Here there is a chapter in dogmatics
which the theologians, in all the many years during which theology has
flourished, have neglected to write, namely: a new or higher logic and theory of
knowledge, which could show us by the help of what new principles one slips
past, in faith, the reefs on which natural thinking constantly runs aground.
Until this chapter is written, we must have the right to examine theological
concepts in the same way as we examine all others, and to pass the same judgment
on inconsistencies and unwarranted applications of concepts as we pass
everywhere else. Such an examination will lead to the result that the dogma
stands in reality just as powerless before the great riddles as science does.

\medskip
\noindent 2. But --- it will then be said from the theological side (especially
now, when philosophical interest in our theological world seems to have been
laid in the grave with Bishop Martensen) --- it is after all a misunderstanding
to suppose that faith is to satisfy a theoretical need. For faith, thought and
reflection mean nothing, and all the counter-arguments in the world fall
powerless to the ground for the simple reason that no one comes to faith except
the one who seeks consolation for an anguished conscience. It is feeling,
particularly the anguished conscience, and no other force, that is at work here.

This is not an agreeable objection, if what is meant is that the speaker himself
cannot do without the consolation that lies in faith; for in that case the
discussion must cease. Who would rob anyone of his consolation? Even the most
zealous chemist would not use for analysis the drink of water with which a human
being is to slake his burning thirst. Every discussion presupposes that the
participants each have so much calm and strength of mind that they can look the
truth in the face, so that one need not, as with the sick, conceal anything.

What consoles the one, moreover, does not console the other, or even arouses his
offense and indignation. As soon as an appeal is made to feeling, the appeal is
more or less to something purely individual. Such an appeal can never be an
argument; for feeling in and for itself teaches us nothing whatever beyond our
own state. The most vivid joy and reassurance that I feel at a piece of news
guarantees nothing about that news's reliability. From my feeling I can
therefore infer nothing with certainty about anything that lies beyond my own
state. Our wishes do not determine what is truth; but it might well be that
H.~C.~Ørsted was right that the full truth always brings its consolation with
it.

The feeling to which appeal is especially made as a support for theological
faith is the consciousness of sin. Here now arises the great question whether
the consciousness of sin, as theology conceives it, really is, ethically
regarded, a sound and warranted feeling --- whether there may not be contained
in it elements that are not really ethical, but hang together with barbaric
representations (of an angry deity who must be propitiated), and which must be
separated out if the consciousness of sin is to become a real ethical feeling of
remorse.

But it is precisely the measure of the ethical that is so different at the two
standpoints! And theology, which formerly referred us from the dogmatic over to
the ethical, now maintains conversely that without dogmatic faith no ethics is
possible. The appeal to the consciousness of sin therefore leads us only in a
circle: for the consciousness of sin is a feeling conditioned by definite
dogmatic representations, and it is then no wonder if it bears witness in favor
of the dogma. --- We get out of this circle only if theology can show us --- what
has hitherto not succeeded --- that only dogmatic faith can be the foundation of
ethics; indeed, that the theological assumptions are on the whole suited to be
the foundation of an ethics. If, on the other hand, a humane ethics is possible,
and if an unconscious humane ethics should prove always to underlie theological
ethics in its various forms (as one might infer from the fact that when human
beings improve, their gods improve too), then the sharpest sting would be taken
away, and what then remained of the problem might well continue to set feeling
in strong and deep motion, but would no longer split our race into two opposed
camps that clash in irreconcilable strife over the ultimate measure of good and
evil.

\medskip
\noindent 3. There is an assertion that has long been waiting impatiently to be
heard from the theological side. It runs roughly thus: ``Faith is built upon
historical facts; its content is first and foremost history, and neither some
excogitated system nor mere postulates of feeling. And historical facts one must
take as they are; no argument of reason can shake them.''

If only the word history were here taken in the same sense in which we take it in
historical science! But it goes with the theological use of this concept as with
the use of other concepts (cause, ethics, and so on): despite the identity of the
word, the real meaning of the concepts is very different. Historical is not
everything that stands in old books or is contained in old traditions. A
historical fact may often require penetrating critical inquiry before it can lay
claim to acknowledgment. The great question then becomes whether theology will
submit to the ordinary laws of historical criticism, or whether it lays claim to
certain privileges for its ``facts.'' That it must demand such privileges follows
already from the fact that its source writings contain narratives of miracles
that require faith in order to be acknowledged. A historical inquiry that applies
no other presuppositions than those from which all science must set out leads ---
as experience has always shown --- to a wholly different conception of the events
at the founding of the Church from the one that theological faith lays down. It
goes, then, with the appeal to history as with the appeal to feeling and to
ethics: it leads in a circle, for by history is meant history theologically
conceived and evaluated. ---

\end{document}
